% Wersja źródłowa: audyt II, jednostka 19
\setcounter{section}{18}
\section{Przestrzeń dystrybucji i jej topologie}
\label{sec:dystrybucje19}

\providecommand{\atlasref}[2]{#2}


\subsection{Definicja i pierwsze przykłady}

\noindent\emph{Powiązanie tematyczne w Bibliotece matematycznej:} \atlasref{PDE-07::def:pde7-distribution}{PDE-07: dystrybucje w teorii równań różniczkowych}.


W całym rozdziale zakładamy, że $d\ge1$ jest liczbą całkowitą,
$\Omega\subseteq\R^d$ jest niepustym zbiorem otwartym,
a $\K\in\{\R,\C\}$.
Przestrzeń $\mathcal D(\Omega;\K)$ ma topologię skonstruowaną
w \hyperref[sec:funkcje-testowe18]{rozdziale~18}. W szczególności
jest lokalnie wypukła, Hausdorffa i zupełna, ale nie jest
metryzowalna. Dla $K\Subset\Omega$ zapisujemy
$\mathcal D_K=\{\varphi\in\mathcal D(\Omega):\supp\varphi\subseteq K\}$
oraz $p_m(\varphi)=\max_{|\alpha|\le m}\sup_{x\in\Omega}
|\partial^\alpha\varphi(x)|$. Symbolem $\mathcal D'(\Omega;\K)$
oznaczamy zbiór \emph{ciągłych $\K$-liniowych} funkcjonałów
na $\mathcal D(\Omega;\K)$.

\begin{definicja}[Dystrybucja]
\label{def:distribution19}
\emph{Dystrybucją} na $\Omega$ nazywamy element
$T\in\mathcal D'(\Omega;\K)$; wartość $T$ na funkcji
testowej $\varphi$ zapisujemy $\langle T,\varphi\rangle$.
Przyjmujemy konwencję $\K$-liniowości względem $\varphi$:
$\langle T,a\varphi+b\psi\rangle=a\langle T,\varphi\rangle+
b\langle T,\psi\rangle$ dla dowolnych $a,b\in\K$
i $\varphi,\psi\in\mathcal D(\Omega;\K)$.
Nie wprowadzamy tutaj sprzężenia
zespolonego w drugim argumencie.
\end{definicja}

\begin{twierdzenie}[Lokalne oszacowanie charakteryzujące dystrybucje]
\label{thm:distribution-estimate19}
Niech $T:\mathcal D(\Omega;\K)\to\K$ będzie odwzorowaniem
liniowym. Następujące warunki są równoważne:
\begin{enumerate}
\item $T$ jest ciągły w standardowej topologii $\mathcal D(\Omega)$;
\item dla każdego $K\Subset\Omega$ istnieją $C_K\ge0$
i $m_K\in\N_0$ takie, że
\[|\langle T,\varphi\rangle|\le C_K p_{m_K}(\varphi)
\qquad(\varphi\in\mathcal D_K(\Omega)).\]
\end{enumerate}
\end{twierdzenie}
\begin{proof}
Z własności uniwersalnej granicy induktywnej
(Twierdzenie~\ref{thm:LF-universal18}) ciągłość $T$
jest równoważna ciągłości każdego
$T|_{\mathcal D_K}$. Topologia $\mathcal D_K$ jest wyznaczona
przez rosnącą rodzinę seminorm $(p_m)_{m\in\N_0}$.
Jeśli $T|_{\mathcal D_K}$ jest ciągły, istnieje $m$ oraz
$\varepsilon>0$ takie, że $p_m(\varphi)<\varepsilon$
implikuje $|T(\varphi)|<1$. Dla $p_m(\varphi)>0$ stosujemy
tę implikację do $a\varphi$ dla
$a=\varepsilon/(2p_m(\varphi))>0$.
Wówczas $p_m(a\varphi)=\varepsilon/2$ oraz
$a|T(\varphi)|<1$, czyli
$|T(\varphi)|\le(2/\varepsilon)p_m(\varphi)$.
Jeśli $p_m(\varphi)=0$, to dla każdego $a>0$ mamy
$a|T(\varphi)|<1$, skąd $T(\varphi)=0$. Odwrotnie, podane oszacowanie
zapewnia ciągłość na każdym etapie, a więc na całej
$\mathcal D(\Omega)$.
\end{proof}

\begin{definicja}[Rząd dystrybucji na zwartym zbiorze]
\label{def:order19}
Dla $T\in\mathcal D'(\Omega;\K)$, zwartego $K\Subset\Omega$
i $m\in\N_0$ mówimy,
że $T$ ma \emph{rząd nie większy niż $m$ na $K$}, jeżeli
istnieje $C_K\ge0$ takie, że
$|\langle T,\varphi\rangle|\le C_K p_m(\varphi)$
dla $\varphi\in\mathcal D_K$. Mówimy, że $T$ ma
\emph{rząd skończony na $\Omega$}, jeżeli można dobrać
jedno $m$ właściwe dla wszystkich $K\Subset\Omega$
(przy czym $C_K$ może zależeć od $K$).
Dokładnym rzędem na $K$ nazywamy najmniejszą liczbę
$m\in\N_0$, dla której istnieje $C_K\ge0$ spełniająca
powyższe oszacowanie dla wszystkich $\varphi\in\mathcal D_K$.
Istnienie takiego minimum wynika z
Twierdzenia~\ref{thm:distribution-estimate19} i dobrego
uporządkowania $\N_0$.
\end{definicja}

\begin{uwaga}[Zależność od zbioru zwartego]
\label{rem:order-local19}
Twierdzenie~\ref{thm:distribution-estimate19} gwarantuje
\emph{lokalnie} skończony rząd każdej dystrybucji,
nie zaś istnienie jednego rzędu na całej $\Omega$.
Przykład~\ref{ex:variable-order18} konstruuje dystrybucję
$\langle T,\varphi\rangle=\sum_{n\ge1}\varphi^{(n)}(n)$
na $\R$, która nie ma skończonego rzędu globalnego.
\end{uwaga}

\begin{przyklad}[Dystrybucje regularne]
\label{ex:regular19}
Jeżeli $f\in L^1_{\mathrm{loc}}(\Omega;\K)$, to
\[\langle T_f,\varphi\rangle=\int_\Omega
f(x)\varphi(x)\,\mathrm d\lambda_d(x)\]
określa funkcjonał liniowy, niezależny od wyboru reprezentanta
klasy $f$. Dla $\varphi\in\mathcal D_K$ całka jest bezwzględnie
zbieżna i zachodzi
\[|\langle T_f,\varphi\rangle|\le
\left(\int_K|f(x)|\,\mathrm d\lambda_d(x)\right)p_0(\varphi).\]
Na mocy Twierdzenia~\ref{thm:distribution-estimate19}
jest to dystrybucja. Tak więc dystrybucje regularne
mają rząd nie większy niż $0$.
Zwróćmy uwagę, że nie jest potrzebne założenie
$f\in L^1(\Omega)$; np. $f(x)=e^{|x|^2}$ na $\R^d$
definiuje dystrybucję.
\end{przyklad}

\begin{stwierdzenie}[Wierność utożsamienia funkcji z dystrybucją]
\label{prop:regular-injective19}
Jeśli $f,g\in L^1_{\mathrm{loc}}(\Omega;\K)$ i $T_f=T_g$,
to $f=g$ prawie wszędzie w $\Omega$ względem miary Lebesgue'a.
\end{stwierdzenie}
\begin{proof}
Z liniowości całki $T_f-T_g=T_{f-g}$, więc po zastąpieniu
$f$ przez $f-g$ wystarczy rozważyć przypadek $T_f=0$.
Wybierzmy niezerową, nieujemną funkcję
$\rho_0\in C_c^\infty(B(0,1);\R)$, korzystając
z Lematu~\ref{lem:cutoffs17}, i przedłużmy ją zerem.
Jej całka jest dodatnia, gdyż $\rho_0$ jest ciągła
i dodatnia na pewnym niepustym zbiorze otwartym.
Niech $\rho=\rho_0/\int_{\R^d}\rho_0\,\mathrm d\lambda_d$.
Wówczas $\supp\rho\subseteq\overline{B(0,1)}$ i
$\int_{\R^d}\rho(x)\,\mathrm d\lambda_d(x)=1$. Określmy
$\rho_\varepsilon(x)=\varepsilon^{-d}\rho(x/\varepsilon)$.
Dla $x\in\Omega$ oraz $\varepsilon>0$ takich, że
$\overline B(x,\varepsilon)\subseteq\Omega$, definiujemy
$(f*\rho_\varepsilon)(x)=\int_\Omega f(y)\rho_\varepsilon(x-y)
\,\mathrm d\lambda_d(y)$.
Funkcja $y\mapsto\rho_\varepsilon(x-y)$ należy do
$\mathcal D(\Omega;\K)$, zatem
$(f*\rho_\varepsilon)(x)=0$.

Na każdym $K\Subset\Omega$ funkcje
$f*\rho_\varepsilon$ zbiegają do $f$ w $L^1(K)$.
Dla niepustego $K$ wybierzmy otwarty $V$ taki, że
$K\subseteq V$ i $\overline V\Subset\Omega$.
Można go skonstruować jako skończoną sumę kul o domknięciach
zawartych w $\Omega$, pokrywających $K$.
Przez $g$ oznaczmy przedłużenie $f\chi_V$ zerem na $\R^d$.
Wtedy $g\in L^1(\R^d)$, gdyż $f$ jest całkowalna na
$\overline V$. Odległość $K$ od $\R^d\setminus V$ jest dodatnia.
Jeśli $\varepsilon$ jest od niej mniejsza, to
$\overline B(x,\varepsilon)\subseteq V$ dla wszystkich $x\in K$,
więc sploty $f*\rho_\varepsilon$ i $g*\rho_\varepsilon$
są równe na $K$. Ponadto, po podstawieniu $y=x-\varepsilon z$,
z normalizacji $\rho$ i twierdzenia Tonellego dostajemy
\[\|g*\rho_\varepsilon-g\|_{L^1(\R^d)}\le
\int_{\R^d}\rho(z)\,
\|g(\cdot-\varepsilon z)-g\|_{L^1(\R^d)}\,
\mathrm d\lambda_d(z)\longrightarrow0,\]
co wynika z ciągłości translacji w $L^1$ i twierdzenia
o zbieżności zmajoryzowanej: dla każdego $z$
norma różnicy translacji dąży do zera, a cała funkcja
podcałkowa jest ograniczona przez
$2\|g\|_{L^1}\rho(z)$, która jest całkowalna.
Ponieważ $g=f$ prawie wszędzie na $K$ i
$f*\rho_\varepsilon=0$ na $K$, otrzymujemy
$\int_K|f|\,\mathrm d\lambda_d=0$.
Stosujemy to do przeliczalnego wyczerpania $(K_n)$
z Lematu~\ref{lem:exhaustion17}. Każdy zbiór
$A_n=\{x\in K_n:f(x)\ne0\}$ ma miarę zero,
a $\{x\in\Omega:f(x)\ne0\}=\bigcup_n A_n$ dla ustalonego
reprezentanta $f$. Z przeliczalnej podaddytywności miary
wynika $f=0$ prawie wszędzie w $\Omega$.
\end{proof}

\begin{przyklad}[Dystrybucje punktowe]
\label{ex:delta19}
Dla $a\in\Omega$ i wielowskaźnika $\alpha\in\N_0^d$
funkcjonały
\[\langle\delta_a,\varphi\rangle=\varphi(a),\qquad
\langle T_{a,\alpha},\varphi\rangle
=(-1)^{|\alpha|}\partial^\alpha\varphi(a)\]
są liniowe, a z oszacowania
$|T_{a,\alpha}(\varphi)|\le p_{|\alpha|}(\varphi)$
wynika ich ciągłość (dla $\alpha=0$ otrzymujemy $\delta_a$).
Mają rząd nie większy niż $|\alpha|$.
Żaden z tych funkcjonałów nie jest postaci $T_f$
dla $f\in L^1_{\mathrm{loc}}(\Omega;\K)$.
Istotnie, wybierzmy $\eta\in C_c^\infty(B(0,1);\R)$
równą $1$ w otoczeniu zera (Lemat~\ref{lem:cutoffs17}).
Przedłużmy $\eta$ zerem na $\R^d$ i połóżmy
$\psi(x)=x^\alpha\eta(x)$ dla $x\in\R^d$.
Wtedy $\partial^\alpha\psi(0)=\alpha!\ne0$.
Dla $0<\varepsilon<1$ tak małego, że
$\overline B(a,\varepsilon)\subseteq\Omega$, połóżmy
$\varphi_\varepsilon(x)=\varepsilon^{|\alpha|}
\psi((x-a)/\varepsilon)$. Wtedy
$T_{a,\alpha}(\varphi_\varepsilon)=(-1)^{|\alpha|}\alpha!$
nie zależy od $\varepsilon$, natomiast dla każdego
$f\in L^1_{\mathrm{loc}}(\Omega)$ mamy
$|T_f(\varphi_\varepsilon)|\le
\varepsilon^{|\alpha|}\|\psi\|_\infty
\int_{B(a,\varepsilon)}|f(x)|\,\mathrm d\lambda_d(x)
\to0$.
Ostatnia granica wynika z bezwzględnej ciągłości całki
funkcji $|f|$ na ustalonej zwartej kuli wokół $a$,
ponieważ $\lambda_d(B(a,\varepsilon))\to0$.
Argument obejmuje również $\alpha=0$, czyli $\delta_a$.
Znak wybrano zgodnie z definicją
pochodnej dystrybucyjnej w kolejnym rozdziale:
$T_{a,\alpha}=\partial^\alpha\delta_a$.
\end{przyklad}

\begin{stwierdzenie}[Dokładny rząd pochodnej delty]
\label{prop:delta-order19}
Niech $a\in\Omega$, $\alpha\in\N_0^d$ i niech
$T_{a,\alpha}$ będzie dystrybucją z Przykładu~\ref{ex:delta19}.
Jeśli zwarty $K\Subset\Omega$ zawiera otoczenie punktu $a$,
to $T_{a,\alpha}$ ma na $K$ dokładnie rząd $|\alpha|$.
\end{stwierdzenie}
\begin{proof}
Oszacowanie rzędu $|\alpha|$ wynika z definicji.
Dla $\alpha=0$ rząd nie może być mniejszy niż $0$.
Niech zatem $m\in\N_0$, $m<|\alpha|$.
Wybierzmy $\psi$ jak w Przykładzie~\ref{ex:delta19}.
Dla małego $\varepsilon>0$ funkcja
$\varphi_\varepsilon(x)=\varepsilon^{|\alpha|}
\psi((x-a)/\varepsilon)$ ma nośnik w $K$. Ponadto
$|T_{a,\alpha}(\varphi_\varepsilon)|
=|\partial^\alpha\psi(0)|>0$, podczas gdy
dla każdego $|\beta|\le m$ zachodzi
$\partial^\beta\varphi_\varepsilon(x)=
\varepsilon^{|\alpha|-|\beta|}(\partial^\beta\psi)((x-a)/\varepsilon)$.
Biorąc supremum i maksimum otrzymujemy
$p_m(\varphi_\varepsilon)\le C_m\varepsilon^{|\alpha|-m}
\to0$, gdzie $C_m=\max_{|\beta|\le m}\|\partial^\beta\psi\|_\infty$.
Nie może więc zachodzić oszacowanie rzędu $m$.
\end{proof}

\begin{przyklad}[Wartość główna Cauchy'ego]
\label{ex:principal-value19}
Dla $\varphi\in\mathcal D(\R;\K)$ definiujemy
\[\left\langle\operatorname{pv}\frac1x,\varphi\right\rangle
=\lim_{\varepsilon\downarrow0}
\int_{\{x\in\R:|x|>\varepsilon\}}
\frac{\varphi(x)}{x}\,\mathrm d\lambda_1(x)
=\int_0^\infty\frac{\varphi(t)-\varphi(-t)}{t}\,
\mathrm d\lambda_1(t).\]
Równość wynika z podstawienia $x=-t$ w całce
po $(-\infty,-\varepsilon)$ i przejścia do granicy,
uzasadnionego poniższym oszacowaniem.
Dla $\supp\varphi\subseteq[-R,R]$ oraz $R\ge1$
z tożsamości $\varphi(t)-\varphi(-t)=\int_{-t}^t\varphi'(s)\,\mathrm ds$
(w tym dla funkcji zespolonych) wynika
$|\varphi(t)-\varphi(-t)|\le2t\,p_1(\varphi)$ dla
$0<t\le1$, natomiast dla $1\le t\le R$ licznik
ma moduł nie większy niż $2p_0(\varphi)$. Dlatego
\[\left|\left\langle\operatorname{pv}\frac1x,
\varphi\right\rangle\right|\le
2p_1(\varphi)+2(\log R)p_0(\varphi).\]
W szczególności całka różnicy jest bezwzględnie zbieżna,
a granica istnieje. Dla każdego zwartego $K\subseteq\R$
wybieramy $R\ge1$ z $K\subseteq[-R,R]$ i stosujemy
$p_0\le p_1$. Liniowość wynika z liniowości całki,
a ciągłość z Twierdzenia~\ref{thm:distribution-estimate19}.
Jest to dystrybucja, chociaż $1/x\notin L^1_{\mathrm{loc}}(\R)$.
Symbol $\operatorname{pv}$ oznacza \emph{wartość główną},
a nie zwykłą, bezwzględnie zbieżną całkę.
\end{przyklad}

\subsection{Nośnik i lokalność}

\begin{definicja}[Zerowanie się dystrybucji na zbiorze otwartym]
\label{def:distribution-zero19}
Niech $T\in\mathcal D'(\Omega)$, a $U\subseteq\Omega$
będzie otwarty. Dla $U=\varnothing$ przyjmujemy
$\mathcal D(U;\K)=\{0\}$ i $T|_U=0$.
Dla niepustego $U$ i $\varphi\in\mathcal D(U;\K)$ oznaczmy przez
$e_U\varphi$ jej przedłużenie zerem na $\Omega$.
Jest ono gładkie, gdyż nośnik $\varphi$ jest zwarty
i zawarty w $U$. Na każdym etapie $\mathcal D_K(U)$
odwzorowanie $e_U$ zachowuje seminormy $p_m$ i ma obraz
w $\mathcal D_K(\Omega)$, więc jest ciągłe.
Z Twierdzenia~\ref{thm:LF-universal18} wynika ciągłość
$e_U:\mathcal D(U;\K)\to\mathcal D(\Omega;\K)$.
Definiujemy ograniczenie $T|_U$ przez
$\langle T|_U,\varphi\rangle=\langle T,e_U\varphi\rangle$.
Jest ono dystrybucją jako złożenie ciągłych odwzorowań liniowych.
Piszemy $T|_U=0$, jeśli powyższa wartość jest zerowa
dla każdej $\varphi\in\mathcal D(U;\K)$.
Nośnikiem $T$ jest domknięty względnie zbiór
\[\supp T=\Omega\setminus\bigcup\{U\subseteq\Omega:
U\text{ jest otwarty i }T|_U=0\}.\]
\end{definicja}

\begin{lemat}[Lokalność dystrybucji]
\label{lem:distribution-local19}
Niech $T\in\mathcal D'(\Omega)$ i
$\varphi\in\mathcal D(\Omega)$ spełnia
$\supp\varphi\cap\supp T=\varnothing$. Wtedy
$\langle T,\varphi\rangle=0$.
\end{lemat}
\begin{proof}
Jeśli $\varphi=0$, teza jest oczywista.
Niech $L=\supp\varphi\ne\varnothing$.
Każdy $x\in L$ ma otoczenie otwarte $U_x\subseteq\Omega$,
na którym $T$ się zeruje. Wybierzmy kulę $B_x$
o środku $x$ i dodatnim promieniu, z $\overline B_x\subseteq U_x$.
Ze zwartości $L$ wybieramy skończone pokrycie
$B_{x_1},\ldots,B_{x_N}$.
Z Lematu~\ref{lem:cutoffs17} wybieramy nieujemne
$\eta_j\in C_c^\infty(U_{x_j};\R)$ równe $1$
w otoczeniu $\overline B_{x_j}$ i przedłużamy je zerem.
Niech $h=\sum_{j=1}^N\eta_j$ i $O=\{x\in\Omega:h(x)>0\}$.
Zbiór $O$ jest otwarty i zawiera $L$.
Na $O$ definiujemy $\varphi_j=\eta_j\varphi/h$,
a na $\Omega\setminus O$ przyjmujemy $\varphi_j=0$.
Przedłużenie jest gładkie, ponieważ każdy punkt spoza $O$
ma otoczenie rozłączne z domkniętym zbiorem $L$,
na którym $\varphi_j$ jest identycznie zerowa.
Mamy $\supp\varphi_j\subseteq L\cap\supp\eta_j\Subset U_{x_j}$
oraz $\sum_j\varphi_j=\varphi$ na $O$ i poza $O$.
Zatem $\langle T,\varphi_j\rangle=0$ dla każdego $j$.
Sumując, otrzymujemy tezę.

W szczególności $T$ zeruje się na otwartym zbiorze
$\Omega\setminus\supp T$, gdyż każda funkcja testowa
na tym zbiorze, przedłużona zerem, ma nośnik rozłączny
z $\supp T$. Jest to więc największy otwarty zbiór
zerowania się $T$.
\end{proof}

\begin{stwierdzenie}[Nośniki przykładów]
\label{prop:support-examples19}
Dla $a\in\Omega$ i $\alpha\in\N_0^d$ mamy
$\supp\delta_a=\supp T_{a,\alpha}=\{a\}$.
Dla $f\in L^1_{\mathrm{loc}}(\Omega)$ nośnik $T_f$
jest dopełnieniem największego otwartego podzbioru
$U\subseteq\Omega$, na którym $f=0$ prawie wszędzie.
\end{stwierdzenie}
\begin{proof}
Gdy otwarty $U\subseteq\Omega$ nie zawiera $a$,
przedłużenie każdej funkcji z $\mathcal D(U)$ znika
w otoczeniu $a$, zatem wszystkie jej pochodne w $a$ są zerowe.
Jeśli $a\in U$, wybieramy $\eta\in C_c^\infty(U;\R)$
równą $1$ w otoczeniu $a$. Funkcja
$\varphi(x)=(x-a)^\alpha\eta(x)$ spełnia
$\partial^\alpha\varphi(a)=\alpha!\ne0$.
Stąd żadne otoczenie $a$ nie jest zbiorem zerowania
$T_{a,\alpha}$, także dla $\alpha=0$.
Dla $T_f$ mamy $(T_f)|_U=T_{f|_U}$ z definicji ograniczenia.
Jeśli $f=0$ prawie wszędzie w $U$, to $(T_f)|_U=0$.
Odwrotnie, ze Stwierdzenia~\ref{prop:regular-injective19}
na $U$ wynika, że $(T_f)|_U=0$ pociąga $f=0$ prawie wszędzie w $U$.
Największy taki otwarty zbiór istnieje na mocy Lematu
o lokalności, a jego dopełnienie jest nośnikiem $T_f$.
\end{proof}

\begin{twierdzenie}[Zwarty nośnik implikuje skończony rząd]
\label{thm:compact-support-order19}
Jeżeli $T\in\mathcal D'(\Omega;\K)$ i
$S=\supp T\Subset\Omega$, to istnieją
$m\in\N_0$, $C\ge0$ i $K\Subset\Omega$ takie, że
\[|\langle T,\varphi\rangle|\le C p_{K,m}(\varphi)
\quad(\varphi\in\mathcal D(\Omega)),\qquad
p_{K,m}(\varphi)=\max_{|\alpha|\le m}
\sup_{x\in K}|\partial^\alpha\varphi(x)|.\]
W szczególności $T$ ma skończony rząd na całej $\Omega$.
Ponadto $T$ przedłuża się w sposób jednoznaczny do
ciągłego funkcjonału liniowego na $C^\infty(\Omega;\K)$.
\end{twierdzenie}
\begin{proof}
Jeśli $S=\varnothing$, Lemat~\ref{lem:distribution-local19}
daje $T=0$. Możemy wtedy wybrać $m=0$, $C=0$
i dowolny niepusty zwarty $K\Subset\Omega$.
Przedłużeniem jest funkcjonał zerowy; jego jednoznaczność
wynika z gęstości $\mathcal D(\Omega)$ w $C^\infty(\Omega)$
na mocy Wniosku~\ref{cor:density-cinf17}.
W przeciwnym przypadku z Lematu~\ref{lem:cutoffs17}
wybieramy $\chi\in C_c^\infty(\Omega;\R)$ równą $1$
na pewnym otoczeniu $S$. Wtedy
$\supp((1-\chi)\varphi)\cap S=\varnothing$, więc z
Lematu~\ref{lem:distribution-local19} wynika
$\langle T,\varphi\rangle=\langle T,\chi\varphi\rangle$.
Niech $K=\supp\chi$. Z Twierdzenia~\ref{thm:distribution-estimate19}
istnieją $m,C_0$ takie, że
$|\langle T,\psi\rangle|\le C_0p_m(\psi)$ dla
$\psi\in\mathcal D_K$. Reguła Leibniza daje
$p_m(\chi\varphi)\le 2^m p_m(\chi)p_{K,m}(\varphi)$,
ponieważ wszystkie pochodne $\chi$ znikają poza $K$.
Otrzymujemy żądane oszacowanie z $C=C_0 2^m p_m(\chi)$.
Dla każdego zwartego $L\Subset\Omega$ i $\varphi\in\mathcal D_L$
mamy $p_{K,m}(\varphi)\le p_m(\varphi)$,
więc jeden rząd $m$ wystarcza dla wszystkich $L$.
Dla $u\in C^\infty(\Omega)$ definiujemy
$\widetilde T(u)=\langle T,\chi u\rangle$.
Odwzorowanie to jest liniowe, a powyższe oszacowanie
iloczynu obowiązuje również dla $u\in C^\infty(\Omega)$.
Zatem $|\widetilde T(u)|\le C p_{K,m}(u)$,
co zapewnia ciągłość w topologii $C^\infty(\Omega)$, a wzór z poprzedniego akapitu
pokazuje zgodność z $T$ na funkcjach testowych.
Przestrzeń $\mathcal D(\Omega)$ jest gęsta w
$C^\infty(\Omega)$ w tej topologii (Wniosek~\ref{cor:density-cinf17}),
więc przedłużenie jest jednoznaczne.
\end{proof}

\subsection{Słaba-* i silna topologia na przestrzeni dystrybucji}

\begin{definicja}[Topologia słaba-*]
\label{def:weakstar19}
Na $\mathcal D'(\Omega)$ definiujemy
$\sigma(\mathcal D'(\Omega),\mathcal D(\Omega))$ jako
najsłabszą lokalnie wypukłą topologię, dla której
wszystkie funkcjonały $T\mapsto\langle T,\varphi\rangle$,
$\varphi\in\mathcal D(\Omega)$, są ciągłe.
Jest ona wyznaczona przez seminormy
$q_\varphi(T)=|\langle T,\varphi\rangle|$ dla $\varphi\in\mathcal D(\Omega)$.
Rodzina ta oddziela punkty, gdyż dla $T\ne S$ istnieje
$\varphi$ z $\langle T-S,\varphi\rangle\ne0$.
Topologia słaba-* jest zatem Hausdorffa.
Sieć $(T_i)$ zbiega w tej topologii do $T$ dokładnie wtedy,
gdy $\langle T_i,\varphi\rangle\to\langle T,\varphi\rangle$
dla każdej $\varphi\in\mathcal D(\Omega)$.
\end{definicja}

\begin{definicja}[Topologia silna]
\label{def:strong19}
Na $\mathcal D'(\Omega)$ definiujemy
\emph{silną topologię przestrzeni dualnej}
$\beta(\mathcal D'(\Omega),\mathcal D(\Omega))$
jako topologię zbieżności jednostajnej na zbiorach
ograniczonych $B\subseteq\mathcal D(\Omega)$.
Wyznaczają ją seminormy
\[q_B(T)=\sup_{\varphi\in B}|\langle T,\varphi\rangle|.\]
Dla $B=\varnothing$ przyjmujemy $q_B=0$.
Wartości $q_B$ są skończone, ponieważ ciągły funkcjonał
liniowy przenosi zbiory ograniczone na ograniczone.
Dla $B=\{\varphi\}$ otrzymujemy $q_B=q_\varphi$,
zatem topologia silna jest nie słabsza od słabej-*
i także jest Hausdorffa. Sieć $(T_i)$ zbiega silnie do $T$
wtedy i tylko wtedy, gdy $q_B(T_i-T)\to0$
dla każdego ograniczonego $B\subseteq\mathcal D(\Omega)$.
\end{definicja}

\begin{twierdzenie}[Beczkowość przestrzeni testowej]
\label{thm:D-barrelled19}
Przestrzeń $\mathcal D(\Omega;\K)$ jest beczkowa.
Każda rodzina $\mathcal A\subseteq\mathcal D'(\Omega;\K)$
ograniczona punktowo, tj. spełniająca
\[\sup_{T\in\mathcal A}|\langle T,\varphi\rangle|<\infty
\quad\text{dla każdego }\varphi\in\mathcal D(\Omega),\]
jest równociągła.
\end{twierdzenie}
\begin{proof}
Niech $U\subseteq\mathcal D(\Omega)$ będzie
beczką, tj. zbiorem domkniętym, absolutnie wypukłym
i pochłaniającym. Ustalmy wyczerpanie $(K_n)$ z Lematu~\ref{lem:exhaustion17}.
Zbiór $U\cap\mathcal D_{K_n}$ jest domknięty dzięki ciągłości
włożenia etapu, a absolutna wypukłość jest dziedziczona
przez przecięcie z podprzestrzenią liniową. Dla każdego
$\varphi\in\mathcal D_{K_n}$ istnieje $t>0$ z $\varphi\in tU$,
więc $\varphi/t\in U\cap\mathcal D_{K_n}$.
Jest to zatem beczka w przestrzeni Frécheta $\mathcal D_{K_n}$,
która jest beczkowa na mocy Wniosku~\ref{cor:frechet-barrelled}. Zatem
$U\cap\mathcal D_{K_n}$ jest otoczeniem zera
w $\mathcal D_{K_n}$.
Funkcjonał Minkowskiego $p_U$ jest seminormą, gdyż $U$
jest absolutnie wypukły i pochłaniający.
Na $U\cap\mathcal D_{K_n}$ mamy $p_U\le1$, zatem dla
każdego $\varepsilon>0$ seminorma $p_U$ jest mniejsza niż
$\varepsilon$ na otoczeniu $(\varepsilon/2)(U\cap\mathcal D_{K_n})$.
Jej ograniczenie do każdego etapu jest więc ciągłe. Własność uniwersalna
granicy induktywnej pokazuje, że $p_U$ jest ciągłą
seminormą na $\mathcal D(\Omega)$; zbiór
$\{\varphi:p_U(\varphi)<1\}\subseteq U$ jest otoczeniem
zera. Przestrzeń $\mathcal D(\Omega)$ jest więc beczkowa.
Druga część wynika z Twierdzenia~\ref{thm:banach-steinhaus-lcs}
dla rodziny ciągłych operatorów liniowych o wartościach w $\K$.
\end{proof}

\begin{wniosek}[Ograniczoność i równociągłość rodzin dystrybucji]
\label{cor:bounded-distributions19}
Dla rodziny $\mathcal A\subseteq\mathcal D'(\Omega)$
następujące warunki są równoważne:
\begin{enumerate}
\item $\mathcal A$ jest punktowo ograniczona na $\mathcal D(\Omega)$;
\item $\mathcal A$ jest równociągła na $\mathcal D(\Omega)$;
\item $\mathcal A$ jest ograniczona w silnej topologii
$\beta(\mathcal D',\mathcal D)$.
\end{enumerate}
\end{wniosek}
\begin{proof}
Pierwszy warunek implikuje drugi przez
Twierdzenie~\ref{thm:D-barrelled19}. Jeżeli
$\mathcal A$ jest równociągła, to dla każdego
$\varepsilon>0$ istnieje otoczenie $V$ zera takie,
że $|\langle T,\varphi\rangle|\le\varepsilon$ dla
$T\in\mathcal A$ i $\varphi\in V$.
Każdy zbiór ograniczony $B$ leży w $\lambda V$
dla pewnego $\lambda>0$, więc
$\sup_{T\in\mathcal A}q_B(T)\le\lambda\varepsilon$.
To daje ograniczoność silną. Wreszcie singletony
$\{\varphi\}$ są ograniczone, więc trzeci warunek
implikuje pierwszy.
\end{proof}

\begin{wniosek}[Jednostajne lokalne oszacowania rodziny dystrybucji]
\label{cor:uniform-order19}
Jeżeli rodzina $\mathcal A\subseteq\mathcal D'(\Omega)$ jest
punktowo ograniczona, to dla każdego $K\Subset\Omega$
istnieją $m_K\in\N_0$ i $C_K\ge0$ takie, że
$|\langle T,\varphi\rangle|\le C_Kp_{m_K}(\varphi)$
jednocześnie dla wszystkich $T\in\mathcal A$
i $\varphi\in\mathcal D_K$.
\end{wniosek}
\begin{proof}
Z Twierdzenia~\ref{thm:D-barrelled19} rodzina $\mathcal A$
jest równociągła na $\mathcal D(\Omega)$, a więc jej
ograniczenia do $\mathcal D_K$ są równociągłe.
Istnieje zatem otoczenie zera $U$ w $\mathcal D_K$,
na którym $|\langle T,\varphi\rangle|\le1$ dla wszystkich
$T\in\mathcal A$. Ponieważ seminormy $p_m$ są rosnące,
istnieją $m_K$ i $\varepsilon_K>0$ takie, że
$\{\varphi:p_{m_K}(\varphi)<\varepsilon_K\}\subseteq U$.
Jeśli $p_{m_K}(\varphi)>0$, podstawiamy
$a\varphi$ dla $a=\varepsilon_K/(2p_{m_K}(\varphi))$.
Wówczas $a|\langle T,\varphi\rangle|\le1$
jednocześnie dla wszystkich $T\in\mathcal A$,
co daje tezę z $C_K=2/\varepsilon_K$.
Jeśli $p_{m_K}(\varphi)=0$, to $a\varphi\in U$
dla każdego $a>0$, więc $\langle T,\varphi\rangle=0$
dla wszystkich $T\in\mathcal A$.
\end{proof}

\begin{twierdzenie}[Zbieżność ciągów: słaba-* = silna]
\label{thm:weak-strong-sequences19}
Niech $(T_n)_{n\ge1}\subseteq\mathcal D'(\Omega;\K)$
i $T\in\mathcal D'(\Omega;\K)$. Następujące warunki
są równoważne:
\begin{enumerate}
\item $T_n\to T$ w $\sigma(\mathcal D',\mathcal D)$;
\item $T_n\to T$ w $\beta(\mathcal D',\mathcal D)$.
\end{enumerate}
\end{twierdzenie}
\begin{proof}
Implikacja (2)$\Rightarrow$(1) wynika z definicji.
Zakładamy (1). Rodzina $\{T_n-T:n\ge1\}$ jest
punktowo ograniczona, gdyż dla ustalonej funkcji
testowej ciąg liczb $\langle T_n-T,\varphi\rangle$
jest zbieżny. Z Wniosku~\ref{cor:bounded-distributions19}
rodzina ta jest równociągła. Ustalmy zbiór ograniczony
$B\subseteq\mathcal D(\Omega)$ i $\varepsilon>0$.
Jego domknięcie $K=\overline B^{\,\mathcal D}$ jest zwarte
na mocy Twierdzenia~\ref{thm:D-Montel18}.
Jeśli $B=\varnothing$, to $q_B=0$, więc załóżmy $B\ne\varnothing$.
Z równociągłości wybierzmy otwarte, zrównoważone otoczenie
zera $V$ takie, że
$|\langle T_n-T,v\rangle|<\varepsilon/2$
za każdym razem, gdy $n\ge1$ i $v\in V$.
Zwartość $K$ pozwala wybrać
$\varphi_1,\ldots,\varphi_N\in K$ z
$K\subseteq\bigcup_{j=1}^N(\varphi_j+V)$.
Zbieżność punktowa zapewnia indeks $n_0$ taki, że
$|\langle T_n-T,\varphi_j\rangle|<\varepsilon/2$
dla wszystkich $n\ge n_0$ i $j\le N$.
Dla dowolnego $\varphi\in B$ piszemy
$\varphi=\varphi_j+v$ z $v\in V$, a więc dla $n\ge n_0$ z liniowości i nierówności trójkąta
$|\langle T_n-T,\varphi\rangle|
\le|\langle T_n-T,\varphi_j\rangle|+|\langle T_n-T,v\rangle|
<\varepsilon$.
Biorąc supremum, otrzymujemy $q_B(T_n-T)\le\varepsilon$
dla $n\ge n_0$. Ponieważ $\varepsilon>0$ jest dowolne,
$q_B(T_n-T)\to0$, co kończy dowód.
\end{proof}

\begin{uwaga}[Dlaczego nie wystarczy zbieżność sieci?]
\label{rem:net-distinction19}
Powyższy wynik dotyczy ciągów. Dla ogólnej sieci
zbieżność punktowa nie musi zapewniać punktowej
ograniczoności \emph{całej} rodziny jej wyrazów.
Nie wolno więc mechanicznie zastępować w tym dowodzie
ciągu przez sieć. Poniżej wykażemy, że obie topologie
są istotnie różne.
\end{uwaga}

\begin{twierdzenie}[Ścisła różnica topologii]
\label{thm:weak-strong-different19}
Jeżeli $\Omega\subseteq\R^d$ jest niepustym zbiorem
otwartym, to na $\mathcal D'(\Omega)$ topologia
$\beta(\mathcal D',\mathcal D)$ jest ściśle silniejsza
od $\sigma(\mathcal D',\mathcal D)$. Istnieje sieć
dystrybucji zbieżna słabo-* do zera, która nie
zbiega silnie do zera.
\end{twierdzenie}
\begin{proof}
Wybierzmy $a\in\Omega$ i $r>0$ tak, że
$\overline B(a,r)\Subset\Omega$. Definiujemy
$\chi(x)=\exp(-1/(r^2-|x-a|^2))$ dla $|x-a|<r$
i $\chi(x)=0$ poza kulą.
Gładkość na brzegu wynika z zanikania wszystkich pochodnych
funkcji $t\mapsto e^{-1/t}$ przy $t\downarrow0$,
jak w dowodzie Lematu~\ref{lem:cutoffs17}.
Mamy $\supp\chi=\overline B(a,r)$ i $\chi>0$ w $B(a,r)$.
Dla $j\ge0$ połóżmy
$\psi_j(x)=2^{-j}((x_1-a_1)/r)^j\chi(x)$.
Niech $u(x)=(x_1-a_1)/r$. Dla $0\le k\le j$ mamy
$\partial_1^k u^j=r^{-k}j!/(j-k)!\,u^{j-k}$,
a dla $k>j$ pochodna jest zerowa.
Na nośniku $\chi$ zachodzi $|u|\le1$.
Każda pochodna $u^j$ zawierająca różniczkowanie względem
którejkolwiek z pozostałych współrzędnych jest zerowa.
Z reguły Leibniza i ograniczoności pochodnych $\chi$
wynika więc $p_m(\psi_j)\le C_m 2^{-j}(1+j)^m$,
gdzie $C_m$ nie zależy od $j$. Ciąg po prawej stronie
jest ograniczony dla każdego $m$, gdyż wykładniczy
spadek dominuje dowolną ustaloną potęgę.
Ponadto $\supp\psi_j\subseteq\overline B(a,r)$,
więc z Twierdzenia~\ref{thm:D-bounded18}
zbiór $B_0=\{\psi_j:j\ge0\}$ jest ograniczony w $\mathcal D$.
Jeżeli $\sum_{j=0}^N c_j\psi_j=0$, to na kuli dzielimy
przez dodatnią $\chi$ i ograniczamy się do $x=a+rt e_1$,
$-1<t<1$. Otrzymujemy wielomian
$\sum_{j=0}^N c_j2^{-j}t^j=0$ dla każdego $t\in(-1,1)$.
Wszystkie $c_j$ są zerowe, więc funkcje są liniowo niezależne.

Niech $I$ będzie zbiorem par $(F,n)$, gdzie $F$ jest
skończonym podzbiorem $\mathcal D(\Omega)$, a $n\in\N$,
z porządkiem $(F,n)\preceq(G,m)$ wtedy i tylko wtedy,
gdy $F\subseteq G$ i $n\le m$. Jest to zbiór skierowany:
wspólnym ograniczeniem górnym par $(F,n)$ i $(G,m)$
jest $(F\cup G,\max\{n,m\})$.
Dla każdej pary wybieramy $\psi\in B_0\setminus\operatorname{lin}(F)$,
co jest możliwe, gdyż $B_0$ zawiera nieskończenie wiele
liniowo niezależnych funkcji, a $\operatorname{lin}(F)$
jest skończenie wymiarowa.
Skończenie wymiarowa podprzestrzeń
$\operatorname{lin}(F)$ jest domknięta, zatem
z Twierdzenia~\ref{thm:point-closed-subspace}
istnieje ciągły funkcjonał liniowy $L$ z
$L|_{\operatorname{lin}(F)}=0$ i $L(\psi)\ne0$.
Definiujemy $T_{F,n}=(n/L(\psi))L$. Wtedy
$T_{F,n}|_{\operatorname{lin}(F)}=0$ i
$\langle T_{F,n},\psi\rangle=n$.

Dla ustalonej $\varphi\in\mathcal D(\Omega)$ i
$(F,n)\succeq(\{\varphi\},1)$ mamy
$\langle T_{F,n},\varphi\rangle=0$. Sieć
$(T_{F,n})_{(F,n)\in I}$ zbiega więc słabo-* do zera.
Jednak $q_{B_0}(T_{F,n})\ge n$ dla każdej pary.
Dla każdego $M>0$, wybierając liczbę całkowitą $n_0>M$,
otrzymujemy $q_{B_0}(T_{F,n})>M$ dla
$(F,n)\succeq(\varnothing,n_0)$. Seminorma dąży więc
do nieskończoności, a sieć nie zbiega silnie do zera. Ponieważ topologia silna
jest nie słabsza od słabej-*, a obie mają różne sieci zbieżne,
jest od niej ściśle silniejsza.
\end{proof}

\subsection{Granice regularizacji i zastosowania}

\begin{przyklad}[Delty uciekające do nieskończoności]
\label{ex:delta-escape19}
Niech $\Omega=\R^d$, a $(a_n)\subseteq\R^d$
spełnia $|a_n|\to\infty$. Wtedy $\delta_{a_n}\to0$
\emph{silnie} w $\mathcal D'(\R^d)$.
Istotnie, każdy ograniczony $B\subseteq\mathcal D(\R^d)$
ma wspólny zwarty nośnik $K$, zatem dla dużych $n$
$\sup_{\varphi\in B}|\varphi(a_n)|=0$.
\end{przyklad}

\begin{przyklad}[Ciągłość delty względem położenia]
\label{ex:delta-continuity19}
Jeżeli $a_n\to a$ w $\Omega$, to
$\delta_{a_n}\to\delta_a$ silnie w $\mathcal D'(\Omega)$.
Dla ograniczonego $B\subseteq\mathcal D(\Omega)$
i $n$ dostatecznie dużych odcinek $[a,a_n]$ leży
w $\Omega$, ponieważ $a_n$ od pewnego miejsca
leży w kuli o środku $a$ zawartej w $\Omega$.
Dla każdej $\varphi\in B$ stosujemy tożsamość
$\varphi(a_n)-\varphi(a)=\int_0^1
\sum_{j=1}^d(a_{n,j}-a_j)\partial_j\varphi(a+t(a_n-a))\,
\mathrm dt$.
Nierówność Cauchy'ego--Schwarza obowiązuje również
dla zespolonych wartości pochodnych. Biorąc supremum po $B$,
otrzymujemy
\[q_B(\delta_{a_n}-\delta_a)\le
|a_n-a|\sqrt d\sup_{\varphi\in B}p_1(\varphi),
\]
gdyż $|\nabla\varphi(x)|\le\sqrt d\,p_1(\varphi)$.
Prawa strona dąży do zera na mocy kryterium
ograniczoności z Twierdzenia~\ref{thm:D-bounded18}.
\end{przyklad}

\begin{twierdzenie}[Aproksymacja jednostki i zbieżność do delty]
\label{thm:mollifier-delta19}
Niech $a\in\Omega$ i niech
$\rho\in C_c^\infty(\R^d)$ ma nośnik w
$\overline{B(0,1)}$, spełnia $\rho\ge0$ i
$\int_{\R^d}\rho(x)\,\mathrm d\lambda_d(x)=1$.
Dla $n\ge1$ określmy
$f_n(x)=n^d\rho(n(x-a))$ dla $x\in\Omega$.
Wtedy $T_{f_n}\to\delta_a$ silnie w $\mathcal D'(\Omega)$.
Co więcej, istnieje stała $C_\rho\ge0$ zależna tylko od
$\rho$ i $d$, taka że dla każdego ograniczonego zbioru
$B\subseteq\mathcal D(\Omega)$ i wszystkich dostatecznie
dużych $n$ (niezależnie od $B$) zachodzi
\[q_B(T_{f_n}-\delta_a)\le
\frac{C_\rho}{n}\sup_{\varphi\in B}p_1(\varphi).\]
\end{twierdzenie}
\begin{proof}
Wybieramy $r>0$ z $\overline B(a,r)\subseteq\Omega$,
a następnie $n_0$ tak, że $1/n_0<r$.
Dla $n\ge n_0$ podstawienie $y=n(x-a)$ daje
\[\langle T_{f_n},\varphi\rangle
=\int_{\overline B(0,1)}\rho(y)\varphi(a+y/n)\,
\mathrm d\lambda_d(y).\]
Całkujemy po kuli, gdyż poza nią $\rho=0$,
a dla $y$ w tej kuli argument $\varphi$ należy do $\Omega$.
Dla $y\in\supp\rho$ oraz dostatecznie dużego $n$
odcinek od $a$ do $a+y/n$ leży w $\Omega$. Zatem
z twierdzenia Newtona--Leibniza i nierówności Cauchy'ego--Schwarza
otrzymujemy
$|\varphi(a+y/n)-\varphi(a)|\le
\sqrt d\,|y|\,p_1(\varphi)/n$.
Ponieważ $\int\rho\,\mathrm d\lambda_d=1$,
po odjęciu $\varphi(a)$ całkujemy tę różnicę
z nieujemną wagą $\rho$.
Stąd teza z
$C_\rho=\sqrt d\int_{\R^d}|y|\rho(y)\,
\mathrm d\lambda_d(y)$. Ograniczoność $B$ zapewnia
$\sup_{\varphi\in B}p_1(\varphi)<\infty$, więc prawa
strona dąży do zera.
\end{proof}



\begin{twierdzenie}[Zbieżność w $L^1_{\mathrm{loc}}$ a zbieżność dystrybucji]
\label{thm:L1-to-strong19}
Niech $f_n,f\in L^1_{\mathrm{loc}}(\Omega;\K)$ i załóżmy,
że $f_n\to f$ w $L^1(K)$ dla każdego $K\Subset\Omega$.
Wówczas $T_{f_n}\to T_f$ w silnej topologii
$\beta(\mathcal D'(\Omega),\mathcal D(\Omega))$.
\end{twierdzenie}
\begin{proof}
Niech $B\subseteq\mathcal D(\Omega)$ będzie ograniczony.
Jeśli $B=\varnothing$, teza jest natychmiastowa.
Z Twierdzenia~\ref{thm:D-bounded18} istnieje
$K\Subset\Omega$ zawierający wszystkie nośniki funkcji
z $B$, przy czym $M_B=\sup_{\varphi\in B}p_0(\varphi)<\infty$.
Otrzymujemy
\[q_B(T_{f_n}-T_f)\le
M_B\int_K|f_n(x)-f(x)|\,\mathrm d\lambda_d(x)
\longrightarrow0.\]
Jest to dokładnie definicja zbieżności silnej.
\end{proof}

\begin{przyklad}[Silna zbieżność bez zbieżności lokalnej w $L^1$]
\label{ex:oscillations19}
Na $\Omega=\R$ weźmy $f_n(x)=\sin(nx)$.
Dla każdej $\varphi\in\mathcal D(\R)$ wybieramy przedział
$[-R,R]$ zawierający jej nośnik we wnętrzu.
Całkowanie przez części, z zerowymi wyrazami brzegowymi,
daje
\[\langle T_{f_n},\varphi\rangle=
\frac1n\int_\R\cos(nx)\varphi'(x)\,
\mathrm d\lambda_1(x).\]
Niech $B\subseteq\mathcal D(\R)$ będzie ograniczony.
Wspólny nośnik funkcji z $B$ mieści się w $[-R,R]$
dla pewnego $R>0$, a $\sup_{\varphi\in B}p_1(\varphi)<\infty$.
Dlatego
\[q_B(T_{f_n})\le\frac{2R}{n}
\sup_{\varphi\in B}p_1(\varphi)\longrightarrow0.\]
Zatem $T_{f_n}\to0$ \emph{silnie} w $\mathcal D'(\R)$,
ale $\int_{-1}^1|\sin(nx)|\,\mathrm d\lambda_1(x)$
nie dąży do zera. Istotnie, przez parzystość $|\sin|$
i podstawienie $t=nx$ otrzymujemy
$\int_{-1}^1|\sin(nx)|\,\mathrm d\lambda_1(x)
=(2/n)\int_0^n|\sin t|\,\mathrm d\lambda_1(t)$.
Zapisując $n=k\pi+r$ z $k=\lfloor n/\pi\rfloor$ i $0\le r<\pi$,
mamy $\int_0^n|\sin t|\,\mathrm d\lambda_1(t)=2k+\int_0^r\sin t\,\mathrm d\lambda_1(t)$.
Ostatnia całka leży w $[0,2]$, zaś $k/n\to1/\pi$,
zatem granica wynosi $4/\pi$.
Zbieżność dystrybucji jest więc znacznie słabsza niż
zbieżność reprezentujących je funkcji w $L^1_{\mathrm{loc}}$.
\end{przyklad}
